\documentclass[11pt,a4paper]{article}
\usepackage[margin=1in]{geometry}
\usepackage[T1]{fontenc}
\usepackage{lmodern}
\usepackage{array}
\usepackage{booktabs}
\usepackage{enumitem}
\usepackage{xcolor}
\usepackage{hyperref}
\usepackage{fancyvrb}
\hypersetup{colorlinks=true,linkcolor=blue,urlcolor=blue}

\title{Lab 11 Test Report: Symbol Table Manager}
\author{Abdul Rehman}
\date{May 2026}

\begin{document}
\maketitle

\section*{Objective}
This lab implemented a hash-based symbol table manager for a Pascal subset. The main goals were to support identifier insertion, lookup, deletion, nested scopes, duplicate declaration detection, undeclared name detection, and parser integration.

\section*{Implementation Summary}
The symbol table is implemented in \texttt{symbol\_table.py} using separate chaining with a prime table size of 211. Each scope is represented as an independent hash table with a parent pointer. The parser integration is connected to the recursive descent parser from Lab 09, while the semantic driver \texttt{pascal\_semantic\_simulator.py} is used for focused testing.

\section*{Test Cases and Results}
\begin{enumerate}[leftmargin=*, itemsep=0.5em]
    \item \textbf{Task 1: single-scope operations}
    \begin{itemize}
        \item 10 inserts, 5 lookups, and 3 deletes were executed successfully using \texttt{task1\_driver.py}.
        \item Lookups included misses, and deletes correctly removed entries from the current scope.
    \end{itemize}

    \item \textbf{Valid nested scope program}
\begin{Verbatim}[fontsize=\small]
python lab_11/pascal_semantic_simulator.py lab_11/test_cases/valid_nested.pas
\end{Verbatim}
Expected: no semantic errors.
Actual: passes with nested scope dumps, and all lookups resolve correctly.

    \item \textbf{Duplicate declaration}
\begin{Verbatim}[fontsize=\small]
python lab_11/pascal_semantic_simulator.py lab_11/test_cases/duplicate_decl.pas
\end{Verbatim}
Expected: duplicate declaration error for \texttt{x}.
Actual:
\begin{Verbatim}[fontsize=\small]
ERROR line 4: duplicate declaration 'x'
\end{Verbatim}

    \item \textbf{Undeclared variable use}
\begin{Verbatim}[fontsize=\small]
python lab_11/pascal_semantic_simulator.py lab_11/test_cases/undeclared_use.pas
\end{Verbatim}
Expected: undeclared variable error for \texttt{a}.
Actual:
\begin{Verbatim}[fontsize=\small]
ERROR line 5: undeclared variable 'a'
\end{Verbatim}

    \item \textbf{Shadowing in nested scope}
\begin{Verbatim}[fontsize=\small]
python lab_11/pascal_semantic_simulator.py lab_11/test_cases/shadowing_valid.pas
\end{Verbatim}
Expected: inner \texttt{x} should shadow outer \texttt{x} inside the nested block.
Actual: the inner declaration is used inside scope 2, and the outer declaration is restored after scope exit.
\end{enumerate}

\section*{Observed Output Summary}
The output from the tested programs showed the following behavior:
\begin{itemize}
    \item Scope entry and exit messages were printed correctly.
    \item Each scope was dumped in a formatted table at exit.
    \item Duplicate declaration and undeclared variable errors were reported with line numbers.
    \item Shadowing worked as expected, with lookup returning the nearest declaration.
\end{itemize}

\section*{Conclusion}
The lab successfully demonstrates the use of a symbol table in a compiler front-end. The implementation meets the required operations for single-scope hashing, nested scope management, error detection for duplicate and undeclared names, and debug printing. The recursive descent parser integration provides a realistic path for semantic analysis in the final compiler project.

\section*{Notes for Viva}
Be ready to explain:
\begin{itemize}
    \item why a hash table with chaining was chosen,
    \item the difference between \texttt{lookup\_current} and \texttt{lookup},
    \item how shadowing works,
    \item where \texttt{begin\_scope} and \texttt{end\_scope} are called,
    \item and why the table size is a prime number.
\end{itemize}

\end{document}
